\subsection{Overview}

The Work Breakdown Structure (WBS) is a hierarchical decomposition of the project into manageable tasks and deliverables. It provides a structured roadmap for planning, monitoring, and controlling the development process. Each major phase is divided into smaller work packages that can be scheduled, implemented, and evaluated independently.

The WBS for Dikaios Amanah follows the Software Development Life Cycle (SDLC) and ensures that every aspect of the project is systematically planned.

%----------------------------------------------------

\subsection{Level 1 Work Breakdown Structure}

\begin{table}[H]
\centering
\begin{tabular}{|c|p{10cm}|}
\hline
\textbf{WBS ID} & \textbf{Major Phase} \\
\hline
1.0 & Project Initiation \\
\hline
2.0 & Requirements Analysis \\
\hline
3.0 & System Design \\
\hline
4.0 & Database Design \\
\hline
5.0 & Frontend Development \\
\hline
6.0 & Backend Development \\
\hline
7.0 & Security Implementation \\
\hline
8.0 & Trust and Audit Layer Development \\
\hline
9.0 & Testing and Quality Assurance \\
\hline
10.0 & Documentation \\
\hline
11.0 & Final Deployment and Demonstration \\
\hline
\end{tabular}
\caption{Level 1 Work Breakdown Structure}
\end{table}

%----------------------------------------------------

\subsection{Detailed Work Breakdown Structure}

\begin{longtable}{|p{2cm}|p{11cm}|}
\hline
\textbf{WBS ID} & \textbf{Task Description} \\
\hline
1.1 & Literature Review \\
\hline
1.2 & Problem Identification \\
\hline
1.3 & Project Proposal Preparation \\
\hline
1.4 & Requirement Gathering \\
\hline
2.1 & Functional Requirement Analysis \\
\hline
2.2 & Non-functional Requirement Analysis \\
\hline
2.3 & Stakeholder Identification \\
\hline
2.4 & Use Case Analysis \\
\hline
3.1 & System Architecture Design \\
\hline
3.2 & UI/UX Wireframe Design \\
\hline
3.3 & UML Diagram Preparation \\
\hline
3.4 & API Design \\
\hline
4.1 & ER Diagram Design \\
\hline
4.2 & Database Schema Design \\
\hline
4.3 & Database Normalization \\
\hline
5.1 & Authentication Pages \\
\hline
5.2 & Business Owner Dashboard \\
\hline
5.3 & Investor Dashboard \\
\hline
5.4 & Administrator Dashboard \\
\hline
5.5 & Investment Campaign Interface \\
\hline
6.1 & Authentication API \\
\hline
6.2 & Business Management API \\
\hline
6.3 & Investment Management API \\
\hline
6.4 & Document Management API \\
\hline
6.5 & Reporting API \\
\hline
7.1 & JWT Authentication \\
\hline
7.2 & Password Hashing (bcrypt) \\
\hline
7.3 & Role-Based Access Control \\
\hline
7.4 & SHA-256 Document Hashing \\
\hline
7.5 & Activity Logging \\
\hline
8.1 & Audit Layer Design \\
\hline
8.2 & Hash Generation Module \\
\hline
8.3 & Immutable Audit Record Storage \\
\hline
9.1 & Unit Testing \\
\hline
9.2 & Integration Testing \\
\hline
9.3 & System Testing \\
\hline
9.4 & Bug Fixing \\
\hline
10.1 & Technical Documentation \\
\hline
10.2 & User Manual \\
\hline
10.3 & Final Report Writing \\
\hline
11.1 & Final Demonstration \\
\hline
11.2 & Project Submission \\
\hline
\end{longtable}

%----------------------------------------------------

\subsection{WBS Summary}

The proposed Work Breakdown Structure divides the project into manageable development phases that align with the Software Development Life Cycle. Each work package represents a clearly defined deliverable, allowing systematic project management, progress monitoring, and resource allocation throughout the development process. The structured approach also minimizes project risks by enabling incremental development and continuous testing.